%! Author = Saeid Yazdani
%! Date = 11/19/2023


\section{Reliability, Safety, and Security}
\label{sec:reliability-safety-security}

The definition of reliability is context-dependent.
In the context of electronics, a system can be considered reliable if it
satisfies its planned operation over the requested period of time.
Intended operation is based on the knowledge of the mission profile: Only if all
environmental conditions are well-defined can we expect proper behavior and a
long lifetime.

The mission profile describes the intended operational parameters, such as
load force and maximum stress for any given period.
We cannot expect the system to be reliable if the mission profile is roughly
misinterpreted.
We could counteract this by designing a safety-critical system and integrated
stress monitors that detect critical anomalies.
There are areas where we experience unwanted influence on a system:

\begin{itemize}[topsep=8pt,itemsep=4pt,partopsep=4pt, parsep=4pt]

    \item \textbf{\emph{Component Production:}} Every system is made from
    components, and any of those components may be the weakest part of a chain.
    Every component must undergo a qualification process to determine lifetime
    and breakdown risk.
    No single component without such a certificate should be used
    in quality designs.

    \item \textbf{\emph{Substrate/Carrier Production:}}
    Components are assembled on a substrate (\ac{PCB} or \ac{IC}).
    Substrates may suffer from hard-to-find failure effects (cf. delamination,
    blistering, surface contamination, via problems, etc.),
    and even high-end production processes cannot guarantee error-free products.
    Extended \ac{PCB} qualification is necessary to produce high-quality
    assemblies (nuclear power, automotive and airborne industries have
    stringent qualification rules and procedures).
    %
    \item \textbf{\emph{System Assembly:}} A system is made from many boards,
    connectors, and wire harnesses.
    Each element may introduce an additional risk of failure.
    %
    \item \textbf{\emph{System Operation:}}
    Insufficient cooling methods cause thermal stress on components as
    systems may run at elevated temperatures for an extended amount of time
    because the designers are not skilled enough.
    Thermal overstress during normal life is still a problem.
    We can prevent over-voltage damage from lightning, supply surges, or
    \ac{ESD} by proper design actions (surge limiters).
    Over-current due to a short somewhere in the wiring, shoot-through, or the
    like is much harder to prevent, and fuses may not be effective if the
    over-current is only slightly above nominal values (in the case of extended
    slight over-current we will heat the fuse carriers such that they lose their
    function). Even if the system is designed to withstand stress factors,
    reduction of total lifetime because of degradation that accumulates
    over time is still a possibility.
\end{itemize}

When we design system protection mechanisms, we cannot consider all possible
failure causes if they happen simultaneously.
Instead, we use the divide-and-conquer strategy, and we should start with
a single-fault approach where we expect only one failure at a time.

For safety-critical systems, we must identify all inter-dependent
failure causes or \ac{CCF}, which is a single failure event that
causes multiple components to fail \autocite{Belland}.
Particular care is necessary to cope with such conditions, and if the safety
level of the application raises the assumption that the failure of a single
component would likely cause total system failure, redundancy would be the
last resort to maintain system availability and safety by limiting the influence
of a single component failure to induce a destructive event \autocite{NUREG}.

If there is a hidden common-cause effect that was not known during the design
stage, we run into a system crash without prior notice.
Such situations may arise if we do not have complete insight into all component
parameters.
Mechanical and passive components are commonly not described in full detail
(e.g., missing datasheet parameters).
Such components, therefore, adversely affect system reliability and need to be
studied thoroughly before being added to the \ac{BOM}.
Therefore, it is vital to quantify the probability of component failure
over time and develop the means to pinpoint risky parts of a system.
Having identified risky partitions, we may introduce merely adequate monitoring
to get early warnings in case of near-fatal failures of the partition.

One approach to keep a system alive is to replace a defective partition
on-the-fly, i.e. conveying \ac{MTTR} close to zero.
On-the-fly repair procedure requires proper tools and repair kits at hand at the time
of failure, and the system will become larger in size, complexity, and number
of components.
Therefore, we keep such methods at a minimum.

Another substantial factor for system availability is to monitor and detect
early warnings on forthcoming wear-out failures, known as predictive maintenance
so that we replace critical system slices in time. Instead of changing the system
blocks on a fixed timeline (i.e. preventive maintenance), we use hardware to
its near-end operational life, thus reducing downtime and repair or replacement
costs \autocite{womack}.